% ---------------------------------------------------------------------
\subsection{Introduzione}
\label{sec:introduzione}
% ---------------------------------------------------------------------

Questo documento rappresenta la Specifica dei Requisiti Software (SRS) per il sistema \productname, un'applicazione per la gestione collaborativa di issue in progetti software. Il documento è conforme allo standard IEEE 830-1998.

% ---------------------------------------------------------------------
\subsubsection{Scopo}
\label{subsec:scopo}
% ---------------------------------------------------------------------

Lo scopo di questo documento è fornire una specifica completa, accurata e verificabile dei requisiti funzionali e non funzionali del sistema \productname. Questo SRS è destinato ai seguenti stakeholder:

\begin{itemize}[leftmargin=1.1em]
\item \textbf{Team di sviluppo}: per comprendere nel dettaglio le funzionalità da implementare e i vincoli da rispettare.
\item \textbf{Project manager}: per pianificare attività, risorse e tempistiche di sviluppo.
\item \textbf{Quality Assurance}: per definire criteri di accettazione e test plan.
\item \textbf{Committente}: per validare che i requisiti soddisfino le esigenze di business.
\end{itemize}

Il documento costituisce la base contrattuale tra il team di sviluppo e il committente (SoftEngUniNA), definendo in modo non ambiguo ciò che il sistema deve fare e come deve comportarsi.

% ---------------------------------------------------------------------
\subsubsection{Ambito del prodotto}
\label{subsec:ambito}
% ---------------------------------------------------------------------

\noindent\textbf{Nome del prodotto}: \productname

\noindent\textbf{Descrizione generale}: \productname{} è una piattaforma distribuita per la gestione collaborativa di issue (segnalazioni) in progetti software. Il sistema consente di:

\begin{itemize}[leftmargin=1.1em]
\item Organizzare il lavoro in progetti distinti, ciascuno con il proprio team
\item Segnalare issue (bug, feature, domande, documentazione) all'interno dei progetti autorizzati
\item Monitorare lo stato di avanzamento delle issue attraverso workflow configurabili
\item Assegnare issue a membri del team del progetto
\item Collaborare tramite commenti e discussioni
\item Generare report mensili per progetto per analizzare performance del team
\end{itemize}

\noindent\textbf{Gestione Progetti e Permessi}:

Gli amministratori hanno controllo completo su:
\begin{itemize}[leftmargin=1.1em]
\item Creazione, modifica ed eliminazione progetti
\item Assegnazione utenti ai progetti (chi può accedere a cosa)
\item Gestione membri del team per ciascun progetto
\end{itemize}

Gli utenti standard possono:
\begin{itemize}[leftmargin=1.1em]
\item Visualizzare solo i progetti a cui sono stati assegnati dagli amministratori
\item Creare e gestire issue solo all'interno dei progetti a cui hanno accesso
\item Lavorare sulle issue che gli vengono assegnate dagli amministratori
\end{itemize}

\noindent\textbf{Benefici attesi}:

\begin{itemize}[leftmargin=1.1em]
\item Miglioramento della tracciabilità dei problemi e delle richieste
\item Riduzione del tempo medio di risoluzione grazie a una migliore organizzazione
\item Aumento della trasparenza e della collaborazione all'interno del team
\item Supporto decisionale attraverso analisi quantitative delle performance
\end{itemize}

\noindent\textbf{Obiettivi principali}:

\begin{itemize}[leftmargin=1.1em]
\item Fornire un sistema performante (tempo di risposta $\leq$ 3000ms per operazioni comuni)
\item Garantire la sicurezza dei dati sensibili attraverso autenticazione robusta
\item Assicurare l'usabilità con un'interfaccia intuitiva che richieda training minimo
\item Facilitare la manutenibilità attraverso un'architettura distribuita e modulare
\end{itemize}

\noindent\textbf{Cosa NON è incluso nell'ambito}:

\begin{itemize}[leftmargin=1.1em]
\item Integrazione con sistemi di versionamento del codice (Git, SVN)
\item Funzionalità di continuous integration/deployment
\end{itemize}

% ---------------------------------------------------------------------
\subsubsection{Definizioni, acronimi e abbreviazioni}
\label{subsec:definizioni}
% ---------------------------------------------------------------------

Per una trattazione completa dei termini specifici del dominio, degli acronimi e delle abbreviazioni utilizzati in questo documento, si rimanda alla Sezione~\ref{sec:glossario} (Glossario).

Alcuni termini chiave includono:

\begin{itemize}[leftmargin=1.1em]
\item \textbf{Issue}: segnalazione relativa al progetto (bug, feature, question, documentation)
\item \textbf{BE}: Back-end, componente server-side del sistema
\item \textbf{FE}: Front-end, componente client-side del sistema
\item \textbf{API}: Application Programming Interface, interfaccia REST per la comunicazione BE-FE
\end{itemize}

% ---------------------------------------------------------------------
\subsubsection{Riferimenti}
\label{subsec:riferimenti}
% ---------------------------------------------------------------------

Questo documento fa riferimento ai seguenti standard, specifiche e documenti tecnici:

\begin{longtable}{p{0.15\linewidth} p{0.78\linewidth}}
\toprule
\textbf{Rif.} & \textbf{Descrizione} \\
\midrule
\endhead
\bottomrule
\endfoot

RIF-1 & \textit{Progetto-SWE-2025-2026 - canale 2\_v1.pdf}. Specifica del committente con requisiti di alto livello e vincoli del progetto. \\
\addlinespace

RIF-2 & \textit{IEEE Std 830-1998: IEEE Recommended Practice for Software Requirements Specifications}. Standard di riferimento per la struttura e il contenuto di documenti SRS. \\
\addlinespace

RIF-3 & \textit{Cockburn, A. (2000). Writing Effective Use Cases}. Metodologia per la descrizione testuale strutturata di casi d'uso. \\
\addlinespace

\end{longtable}

\newpage

% ---------------------------------------------------------------------
\subsection{Panoramica del documento}
\label{subsec:panoramica}
% ---------------------------------------------------------------------

Il presente documento è organizzato come segue:

\begin{itemize}[leftmargin=1.1em]
\item \textbf{Sezione~\ref{sec:introduzione} -- Introduzione}: fornisce una panoramica del documento, descrive lo scopo del sistema, definisce gli acronimi e elenca i riferimenti normativi e bibliografici.

\item \textbf{Sezione~\ref{sec:glossario} -- Glossario}: contiene le definizioni di tutti i termini tecnici, gli acronimi e i riferimenti utilizzati nel documento per evitare ambiguità.

\item \textbf{Sezione~\ref{sec:casi-uso} -- Casi d'Uso}: presenta il diagramma UML dei casi d'uso e la lista completa delle funzionalità del sistema, modellando le interazioni tra gli attori e il sistema.

\item \textbf{Sezione~\ref{sec:caratterizzazione-utenti} -- Caratterizzazione Utenti}: identifica e descrive in dettaglio i profili degli utenti del sistema (Utente, Utente Autenticato, Amministratore), specificandone obiettivi, responsabilità e permessi.

\item \textbf{Sezione~\ref{sec:requisiti-specifici} -- Requisiti Specifici}: costituisce il cuore tecnico del documento e include:
  \begin{itemize}
  \item \textbf{Requisiti di Sistema}: vincoli architetturali e di sistema imposti dal committente
  \item \textbf{Requisiti Funzionali}: descrizione dettagliata di tutte le funzionalità assegnate
  \item \textbf{Requisiti Non Funzionali}: requisiti di qualità misurabili e verificabili (usabilità, performance, sicurezza, manutenibilità, affidabilità)
  \end{itemize}

\item \textbf{Sezione~\ref{sec:cockburn} -- Descrizione Dettagliata Casi d'Uso}: fornisce la formalizzazione testuale strutturata secondo il template di Cockburn per un caso d'uso significativo.

\item \textbf{Sezione~\ref{sec:mockup} -- Mockup}: presenta la prototipazione visuale delle interfacce utente per i casi d'uso formalizzati, collegata agli step delle tabelle Cockburn.
\end{itemize}

Questo documento deve essere letto in modo sequenziale al primo approccio. Le sezioni successive possono essere consultate in modo non lineare come riferimento durante le fasi di progettazione, implementazione e testing.

\newpage